\documentclass[11pt]{article}
\usepackage[letterpaper,margin=1.05in]{geometry}
\usepackage[T1]{fontenc}
\usepackage{lmodern,microtype,amsmath,amssymb,amsthm,booktabs}
\usepackage[hidelinks]{hyperref}
\newtheorem{theorem}{Theorem}
\setlength{\parindent}{0pt}
\setlength{\parskip}{.55em}
\raggedbottom
\title{\textbf{One-monomino stabilization\\in near-square tatami rectangles}}
\author{1200}
\date{6 November 2026}
\begin{document}
\maketitle
\begin{abstract}
We extract the fixed-gap one-monomino stabilization law from the
generating function of Alhazov, Morita and Iwamoto. The eventual count is
four at gap two and zero at larger even gaps. We identify the last exception
and distinguish a parity-admissible anchor from an integer cutoff.
\end{abstract}

\section{Statement}
Let $T(r,w,m)$ be the number of tilings of a labelled $r$-by-$w$ grid by
dominoes and exactly $m$ unit squares (monominoes), with no four tile corners
meeting at a grid vertex. Reflected or rotated tilings are counted separately
whenever they are distinct. For $m=1$, the area must be odd, so a positive
odd gap $d=w-r$ allows no tiling, and for even $d$ only odd $r$ is admissible.

\begin{theorem}
For $d=2$, $T(5,7,1)=8$ and $T(r,r+2,1)=4$ for every odd $r\geq7$,
with one tiling for each choice of the monomino's corner.
For every even $d\geq4$, $T(d+3,2d+3,1)=4$, and $T(r,r+d,1)=0$
for every odd $r\geq d+5$.
\end{theorem}

For even $d\geq2$, the least admissible stabilization anchor is $d+5$.
The least integer cutoff beyond which the admissible subsequence has
its eventual value is $d+4$. An anchor must itself have the prescribed parity.

\section{Literature input}
For odd $r\geq5$ and $w\geq r$, use Alhazov, Morita and Iwamoto's
one-monomino enumeration~\cite[\S\S3, 5 and 7]{ami}. In the variable $z$
marking width, put $P=z^{r-1}+z^{r+1}$. The count at width $w$ is the
coefficient of $z^w$ in $B+C$, where
\[
 B=\frac{4(z^{r-2}+z^r)}{1-P},
 \qquad C=\frac{2(z+z^r)}{(1-P)^2}.
\]
In $B$, the numerator marks the two special corner block widths. In $C$,
the numerator marks the single-column block $L_r$, which has a monomino at
one end, and the width-$r$ block $M_r$; the two factors $(1-P)^{-1}$ describe
lists of pure domino blocks on either side. We use the widths and formulas
from \S7: the second shift in the printed \S6 recurrence for $B$ and the
inequality introducing \S7 are inconsistent with the block enumeration
and intended longer-width range.

\section{Extracting the fixed-gap coefficients}
Expand $(1-P)^{-1}$ as $\sum_{k\geq0}P^k$ and $(1-P)^{-2}$ as
$\sum_{k\geq0}(k+1)P^k$. A single $P$ contributes exponent $r-1$ or $r+1$;
two or more factors have exponent at least $2(r-1)$.
At $w=r+d$, the coefficient offsets from the two monomials in $B$ are
$d+2,d$; the offsets in $C$ are $r+d-1,d$.

For $d=2$ and odd $r\geq7$, both offsets in $B$ and the second offset in
$C$ are positive and less than $r-1$. The first offset in $C$ is $r+1$,
which comes from one $P$ factor and yields $2\cdot2=4$. These tilings use
$L_r$ and one width-$(r+1)$ block, so their monominoes are at the corners.
The two axis reflections act transitively on corners, giving one tiling
at each. At $r=5$ the same $C$ term contributes four, and $B$ contributes
another four from $d+2=r-1$.

For even $d\geq4$ and odd $r\geq d+5$, both $B$ offsets and the second
$C$ offset are smaller than $r-1$ and positive. The first $C$ offset
satisfies $r+1<r+d-1<2(r-1)$, so it contributes nothing either.
At $r=d+3$, the first $B$ offset is $r-1$, contributing four. The other
$B$ offset and second $C$ offset are below $r-1$; the first $C$ offset
is $2r-4$, strictly between $r+1$ and $2(r-1)$ since $r\geq7$.
This proves both the eventual value and the last exception.

\section{What this says about the conjecture}
The last exceptional admissible size is $d+3$ for every positive even
$d$, so the next admissible size $d+5$ is the least anchor. The cutoff
need only exclude $d+3$, giving $d+4$. For instance, $T(7,11,1)=4$,
and all counts from $r=8$ onward are zero at $d=4$. Even if one anchors
to the value at the cutoff itself, the literal integer minimum here is
$8$, whereas the experimental expression $m+d+4$ in the formulation of
Erickson et al.~\cite[\S4.1, Conjecture 3]{ersw} gives $9$.
Restricting the anchor to admissible sizes gives $9$ and removes this
discrepancy. With $d$ odd there is no admissible subsequence for $m=1$.
The square case is covered separately by Erickson and Schurch~\cite{es}.

This note uses a published enumeration and does not resolve the general
fixed-$m$ conjecture. The companion script checks finite coefficient examples.

\textbf{AI contribution.} GLM5.3F (source project credit supplied by the
author). Literature verification and preparation used GPT Sol Max assistance.

\begin{thebibliography}{9}
\bibitem{ami} A. Alhazov, K. Morita and C. Iwamoto, A note on tatami tilings,
\emph{RIMS K\^oky\^uroku} \textbf{1691} (2010), 1--7.
\href{https://www.kurims.kyoto-u.ac.jp/~kyodo/kokyuroku/contents/pdf/1691-01.pdf}{Kyoto University archive}.
\bibitem{ersw} A. Erickson, F. Ruskey, M. Schurch and J. Woodcock,
Auspicious tatami mat arrangements,
\href{https://arxiv.org/abs/1103.3309}{arXiv:1103.3309} (2011).
\bibitem{es} A. Erickson and M. Schurch, Monomer-dimer tatami tilings of
square regions, \emph{J. Discrete Algorithms} \textbf{16} (2012), 258--269.
\href{https://arxiv.org/abs/1110.5103}{arXiv:1110.5103}.
\end{thebibliography}
\end{document}
